
\section{Replaying a declared evidence vintage}
\label{sim:finite-vintage}
The preceding correction becomes an executable calculation once the inference
request says exactly which information remains active. A document identifier
alone is insufficient. Two documents may repeat one observation, and a new
document may replace an earlier assertion rather than supply a second
independent measurement. The model must retain these relations before it
multiplies likelihoods.

Consider a finite set of persistent company types $h\in H$, with declared prior
masses $p(h)$. Each active observation occupies a named slot and supplies a
likelihood $L_j(h)$. We assume that these observation blocks are conditionally
independent given $h$ in this example. Under that assumption, for a selected
evidence vintage $v$ with active slots $A_v$,
\begin{equation}
 Z_v=\sum_{h\in H}p(h)\prod_{j\in A_v}L_j(h),\qquad
 \pi_v(h)=\frac{p(h)\prod_{j\in A_v}L_j(h)}{Z_v}.
 \label{sim:finite-vintage-posterior}
\end{equation}
The normalizer must be positive. A zero value means that this finite model
assigns zero likelihood to the selected evidence; it does not justify returning
an arbitrary distribution. Missing prior probabilities or likelihoods instead
leave the calculation unresolved. Neither circumstance supplies probabilities
for a family that was specified only as a feasible set.

The selected reference, \texttt{FiniteVintageInference/1}, retains the named
types, exact rational prior, observation-model identity and factorization
assumption. A request also names the actor and information cutoff, the active
observation versions and their replacement ancestry. Each observation retains
its source content, information identity, slot, likelihood vector, supplied
actor-availability time and any superseded version. Its content identifier
changes when any of those fields changes. The semantic information identity
answers a different question: whether two records purport to describe the
same informational contribution. A replacement keeps that identity even when
its likelihood or source changes.

The supplied availability time must be no later than the request's cutoff.
An unknown time remains unresolved; a later time is ineligible for that
request. The reference uses the existing exact fixed-offset timestamp
conversion, including fractional seconds, rather than rounding time to a
calendar day. Its convention excludes leap seconds and unspecified offsets.
These checks establish consistency with declared availability. Establishing
when an actual investor, manager or system could obtain the information still
requires a source-selection procedure.

Replacement selects a new head for an existing slot and preserves the chain
to the earlier record. Retraction removes that slot's likelihood contribution,
equivalently replacing it by one for every type. The request cannot select two
heads for one slot, use a missing predecessor, or treat a different information
identity as a revision of the same observation. Several unresolved candidate
revisions therefore require several conditional requests. Selecting a head
does not establish that it is the authoritative version.

The numerical consequences are immediate in the two-type example. After the
correction, the posterior is $(1/4,3/4)$. Suppose a separately declared,
conditionally independent signal has likelihoods $(1/2,1/4)$. Multiplying
gives unnormalized masses $(1/8,3/16)$, whose sum is $5/16$, and hence
\begin{equation}
 \pi_{\rm corrected+signal}=(2/5,3/5).
\end{equation}
Retracting the corrected observation leaves only the new signal; with the
original equal prior the resulting posterior is $(2/3,1/3)$. Retracting the
only observation restores the prior. A type assigned zero prior mass remains
at zero unless the prior itself is changed explicitly. Restoring support lost
through an erroneous observation is distinct from inventing support absent
from the declared model.

Repackaging the corrected observation under another slot would produce a
different answer. With the equal prior, counting its likelihood twice gives
masses proportional to $(1/16,9/16)$ and posterior $(1/10,9/10)$. The reference
rejects two active observations that declare the same semantic information
identity. An earlier implementation checked distinct slots but omitted this
condition: its arithmetic tests passed while the declared information
contract remained incomplete. The repaired example makes the resulting
overconfidence visible. False distinct identities can still conceal repeated
information, just as a false independence assumption can make a well-formed
calculation substantively wrong.

A checkpoint retains its complete request and computed result. Reuse requires
the same model and prior, actor and exact cutoff, with the old active records
forming an unchanged prefix of the new sequence. New independent records can
then update that prefix. A replacement that changes the prefix instead causes
replay from the preserved prior. This avoids the failed operation that starts
from $(1,0)$ and attempts to recover the excluded type by reweighting it.
Corrupt checkpoint contents or incompatible model and information definitions
are rejected rather than silently relabelled.

In this finite reference the old checkpoint calculation is itself recomputed
before reuse, so the test establishes the posterior relation without claiming
a computational saving. Safe reuse from a larger store would require a proved
sufficient dependency set. Particle checkpoints would additionally retain the
likelihood estimate, proposal and ancestor state required by the extended
target. A vector of posterior weights alone is insufficient.

The reference compares complete masses and normalizers over a finite rational
test grid and explicit contrary cases, including a correction, retraction,
duplicate information, unknown times and a later fractional-second timestamp.
An uninformative observation leaves both types alive. This illustrates a
representation in which uncertainty can persist without enumerating a new
copy of every historical node. It establishes the exact finite calculation
under its declared premises; continuous inference, source completeness and
the suitability of the likelihood for a real company remain further questions.


\section{A correction reaches different decision makers at different times}
\label{sim:actor-vintage-decision}
Replaying the corrected history answers an inference question. It does not
yet say which decision maker could use the correction. Suppose the analyst
learns on 1 February 2026 that the earlier observation was wrong, while the
manager receives that correction on 1 March. On 15 February they possess
different information even though they are studying the same company. A
simulation of the manager's choice must use the manager's information.

Continue the hypothetical fixed-versus-floating example. Under the original
observation the posterior is $(1,0)$; under the correction it is
$(1/4,3/4)$. Suppose an unhedged exposure costs zero for the fixed type and
20 for the floating type. A perfect hedge costs 6 in either case. These are
supplied losses in the same units, and both actions are assumed feasible.
They are enough to define this comparison without inventing an observed
hedging price or an accounting transaction.

Let $p$ denote the probability of the floating type under the particular
actor's available information. The expected losses are therefore
\begin{equation}
  L(\text{unhedged}\mid p)=20p,\qquad
  L(\text{hedged}\mid p)=6.
  \label{sim:actor-hedge-loss}
\end{equation}
The original posterior gives an unhedged loss of zero. The corrected
posterior gives 15, so the hedge has the smaller declared expected loss.
At $p=3/10$ the two losses both equal 6 and both actions remain minimizers.
An implementation must retain that tie unless a separate choice rule is
supplied. The manager on 15 February still uses the old posterior, while
the analyst can already select the hedge under the corrected one.

The selected information procedure takes an actor and an exact cutoff,
selects the available unreplaced observation in each information slot, and
passes that request to the finite evaluator. Adding the future correction
does not change the manager's stored 15 January request or its decision.
Later knowledge can motivate a new decision, but it must not silently
replace the information used to reproduce an earlier one.

Unknown availability cannot be treated as known absence. Two incomparable
available revisions also leave the request unresolved; choosing the newer
file arbitrarily would add an authority rule. The checked examples retain
these distinctions, refuse corrupt observation content and fabricated
posterior results, and leave unknown action losses unresolved. They select
actions from the declared feasible set; they do not execute a hedge or infer
its causal effect. Actual actor access, registry completeness, instrument
terms and feasible financing still need evidence.

